\subsection{Amplitude Carries No Usable Information}

The choice to encode symbols in the time domain is validated directly by the
measurements. Received peak amplitude is pinned: across the 90{,}870 pulses of
the sweep it sits at 41{,}370 of 65{,}535 converter counts with a standard
deviation of 71 counts, a coefficient of variation of 0.17\%, and the 67{,}581
pulses of the single-rate dataset give 42{,}171 counts at the same 0.16\%
spread. On a 3.3~V rail that peak corresponds to roughly 2.1~V, which is where
an MCP6002 output saturates, so the invariance is the amplifier railing rather
than a genuinely constant optical input.

The clearest evidence that the peak is uninformative is that it does not move
even when the signal demonstrably does. Between 28.6 and 171.4~bit/s the peak
high-pass response falls by 41\%, from 1{,}742 to 1{,}033 counts, while over
the same span the mean peak amplitude changes by 64 counts, or 0.16\%. A
quantity that holds to two parts in a thousand while the transition amplitude
driving it changes by nearly half is carrying no information about the channel.
An amplitude-discriminating scheme could not have functioned on this front end,
and no signal-strength observable, for ranging or for anything else, is
available from it without reducing the transimpedance gain.

Pulse-width classes, by contrast, remain well separated.
Fig.~\ref{fig:separability} shows the separation between the zero and one
symbols in pooled standard deviations as a function of data rate. Separation
falls from 39.5 at 28.6~bit/s to 12.5 at 171.4~bit/s, but even at the highest
rate the classes are six times further apart than the two standard deviations
at which distributions begin to overlap appreciably.

\begin{figure}[t]
\centering
\includegraphics[width=\columnwidth]{fig_separability.pdf}
\caption{Separation between the ZERO and ONE symbol classes as a function of
data rate. Classes remain well separated at every rate tested, including rates
at which decoding fails frequently.}
\label{fig:separability}
\end{figure}

\subsection{The Failure Mode Is Omission, Not Confusion}

This separation has a direct consequence: symbol \emph{misclassification} is not
what limits the link. Over 67{,}581 symbols at 28.6~bit/s, no substitution
errors occurred. The entire error budget consisted of 93 pulses (0.14\%) whose
durations fell between the threshold decoder's acceptance windows, and
assigning each to its nearest class recovered all 93.

What limits the link is symbols the receiver never detects at all. At
114.3~bit/s, 0.371\% of symbols went undetected, and 52.5\% of transmissions
were missing at least one. Because a byte spans nine pulses, even this
sub-percent per-symbol rate compounds: it corrupts 3.3\% of bytes, and over the
234 symbols of a 25-character message it leaves only 41.9\% of transmissions
intact. The compounding, not the per-symbol rate, is what makes the link
unreliable, and it is the reason a metric averaged over symbols understates the
problem by more than an order of magnitude.

Crucially, every undetected symbol left a recoverable signature. Of the
transmissions missing exactly one symbol, 100\% exhibited either an anomalous
inter-pulse gap or a merged pulse consistent with a missed falling edge. This
held at every rate tested, not only at 114.3~bit/s. The information required to
reconstruct the lost symbol is present in the captured stream; a per-symbol
classifier simply has no mechanism to use it.

\subsection{Decoder Performance Across Data Rates}

Table~\ref{tab:offline} and Fig.~\ref{fig:results} report the offline
evaluation on held-out messages. The sequence decoder outperforms the threshold
decoder at every rate, by margins ranging from 28 to 61 percentage points of
transmission success. It reduces character error rate at every rate as well, but by an
amount that varies considerably: by half or better at 28.6, 57.1 and
114.3~bit/s, and by 48\% and 33\% at 85.7 and 171.4~bit/s respectively.

\begin{table}[t]
\caption{Offline evaluation, held out by message. Approximately 36 test
transmissions per rate.}
\label{tab:offline}
\centering
\begin{tabular}{@{}lrrrr@{}}
\toprule
& \multicolumn{2}{c}{Transmission success} & \multicolumn{2}{c}{Character error} \\
\cmidrule(lr){2-3}\cmidrule(l){4-5}
Rate (bit/s) & Threshold & Sequence & Threshold & Sequence \\
\midrule
28.6  & 72.2\% & \textbf{100.0\%} & 1.44\% & \textbf{0.00\%} \\
57.1  & 11.1\% & \textbf{72.2\%}  & 6.32\% & \textbf{1.58\%} \\
85.7  & 16.7\% & \textbf{52.8\%}  & 5.75\% & \textbf{3.02\%} \\
114.3 & 36.1\% & \textbf{69.4\%}  & 3.88\% & \textbf{1.72\%} \\
171.4 & 31.4\% & \textbf{60.0\%}  & 4.03\% & \textbf{2.69\%} \\
\bottomrule
\end{tabular}
\end{table}

\begin{figure*}[t]
\centering
\includegraphics[width=\textwidth]{fig_results.pdf}
\caption{Offline evaluation across data rates. (a)~Transmission success rate, with
95\% Wilson score intervals; the dashed line is the fraction of held-out
transmissions that arrived complete or lost a single symbol that left a
recoverable signature, a conservative reference level for a decoder operating
on the detected pulse stream. (b)~Character error rate. The sequence decoder made no
errors at all at 28.6~bit/s.}
\label{fig:results}
\end{figure*}

At 28.6~bit/s the decoder reaches 100\% transmission success with zero character errors,
matching the computed ceiling exactly: every error at that rate is a
dead-zone boundary error, and all are eliminated.

Above that rate the decoder recovers between 42\% and 57\% of the headroom
between the transmissions that arrived fully captured and the ceiling. At
57.1~bit/s, half the held-out transmissions arrived complete and the decoder
raises transmission success to 72.2\% against a ceiling of 88.9\%; at 114.3~bit/s it
reaches 69.4\% from 47.2\% complete, against a ceiling of 94.4\%. These
ceilings are computed on the held-out transmissions alone, so that all three
quantities describe the same messages; evaluated over the full sweep instead,
the same bounds read 93.3\% and 95.8\%, and the mismatch would flatter the
decoder by up to nine points at 85.7~bit/s. The residual gap is genuine and is
discussed in Section~\ref{conclusion}.

The dashed line in Fig.~\ref{fig:results}(a) is not monotonic in rate, and
since that invites a search for a mechanism, it is worth stating exactly what
sets it. Because every single-omission transmission carried a recoverable
signature, the reference level reduces to an identity: it is one minus the
fraction of transmissions that lost \emph{two or more} symbols, and nothing
else enters it. Across the five rates those transmissions number 0, 4, 6, 2 and
5 out of roughly 36.

Their ordering follows per-symbol loss exactly. Measured on the held-out
transmissions alone, loss runs 0.000\%, 0.367\%, 0.495\%, 0.335\% and 0.448\%
(the 0.371\% quoted earlier is the figure over the full sweep at
114.3~bit/s), and ranking the rates by loss reproduces the ceiling ordering
inverted without exception: 114.3~bit/s loses fewer symbols than 85.7~bit/s, so
its ceiling sits higher. Payload length cannot contribute, because the same
twelve messages are held out at every rate. The ceiling is therefore a monotone
function of per-symbol loss, and its non-monotonicity in \emph{rate} is
inherited from symbol loss being itself non-monotonic in rate.

The differences are nonetheless small counts. The step from 6 such
transmissions at 85.7~bit/s to 2 at 114.3~bit/s carries Wilson intervals of
7.9--31.9\% and 1.5--18.1\%, which overlap heavily. The ceiling is best read as
a level of roughly 85--95\% above the lowest rate rather than as a curve whose
shape is resolved.

The confidence intervals in Fig.~\ref{fig:results}(a) bear emphasis. With
approximately 36 transmissions per rate, a point estimate near 70\% carries an
interval of roughly $\pm$15 percentage points. The two decoders are separated
by non-overlapping intervals at 28.6, 57.1, 85.7 and 114.3~bit/s. At
171.4~bit/s they overlap, 18.6--48.0\% against 43.6--74.4\%, so the
28.6-point difference at that rate alone is suggestive rather than established.
The finer structure of the per-rate curves is not resolved at this sample size
and should not be read as a trend.

\subsection{Live Closed-Loop Operation}

Fig.~\ref{fig:live} reports live operation with the decoder in the loop. At
28.6~bit/s over 39 transmissions the sequence decoder achieved 100\%
transmission success with zero character errors, against 69.2\% and 1.54\% for
the threshold
decoder. At 114.3~bit/s over 118 transmissions it achieved 44.9\% against 5.9\%
(a factor of 7.6) and halved character error rate from 5.29\% to 2.61\%,
recovering 46 of the 111 transmissions the threshold decoder lost. No pulses
were reported dropped by the receiver in either run.

\begin{figure}[t]
\centering
\includegraphics[width=\columnwidth]{fig_live.pdf}
\caption{Live closed-loop transmission success at two data rates, with 95\% Wilson
score intervals. The decoder ran on a host computer with the receiver streaming
observations and receiving decoded messages over the same serial link.}
\label{fig:live}
\end{figure}

The live figures at 114.3~bit/s are lower than the corresponding offline
figures, and the cause is payload length rather than any difference in the link
or the decoder. The held-out messages that produced the offline result average
174 symbols; the live beacon transmitted a fixed 234-symbol message. At the
measured per-symbol loss of 0.371\%, the probability that a transmission
arrives with no symbol missing falls from 52.4\% to 41.9\% between those
lengths. Restricting the offline data to messages of 200 symbols or more gives
19.4\% of transmissions fully captured, so the live decoder's 44.9\% is well
above the fraction that arrived intact. This is the payload-length dependence
noted in Section~\ref{system_overview} appearing as a first-order effect, and
it is the reason payload length must be held fixed or reported when comparing
experiments.

\subsection{Residual Error Characteristics}

The decoder's remaining failures at 114.3~bit/s are almost exclusively
single-character substitutions: \texttt{Brovn} for \texttt{Brown},
\texttt{Qudck} for \texttt{Quick}, \texttt{Quick\$Brown} for
\texttt{Quick~Brown}. When a symbol is undetected, gap timing localizes it and
the printable-ASCII constraint narrows the candidates, but does not always
determine the answer uniquely: \texttt{Brovn} is valid ASCII, so the structural
constraint cannot reject it.

This error profile is informative. It indicates that the remaining headroom lies
not in better timing evidence but in additional constraints on the payload
itself, which motivates the language prior identified as the highest-value
future extension.
